\section{Problem formulation}

\subsection{Observations, predictors, and state}

Let $x_t\in\mathcal{X}$ denote the observation at step $t$, and let
$Y_t\in\mathcal{Y}$ denote its random target. The cheap adaptive predictor is
$M_F(x_t;\theta_t)$, while $M_D(x_t)$ denotes the expensive teacher
predictor. The routing action is $a_t\in\{F,D\}$, with prediction
\begin{equation}
  \widehat y_t =
  \begin{cases}
    M_F(x_t;\theta_t), & a_t=F,\\
    M_D(x_t), & a_t=D.
  \end{cases}
\end{equation}
Let $\ell(\widehat y_t,Y_t)$ be a predictive loss. Write
$C_D(S_t,x_t)\geq 0$ for the incremental cost of consulting $M_D$, expressed
in units compatible with the loss; the incremental consultation cost of
choosing $F$ is zero.

The state $S_t=(\theta_t,\mathcal{I}_t)$ includes the cheap predictor's
parameters and the information available before routing the current
observation. The decision uses $(S_t,x_t)$. The information may retain the
available history; no finite-dimensional sufficient state is asserted.
The target is used to define risk, but its availability for adaptation is
not assumed.

If the expensive predictor is queried, let $Z_t=M_D(x_t)$. A general
update map on the full state permits
\begin{equation}
  S_t^+=U(S_t,x_t,Z_t)
  \qquad\text{when }a_t=D.
\end{equation}
This map can adapt $\theta_t$ using the teacher response but may leave it
unchanged. The post-query state $S_t^+$ need not be the complete next-step
state $S_{t+1}$. Neither the update mechanism nor the feedback available
under either action is specified at this stage.

\subsection{Sequential objective and action values}

For notation, consider a finite horizon $T$ and a discount factor
$0\leq\gamma\leq 1$, with $\gamma=1$ giving the undiscounted objective.
A policy $\pi$ selects actions using the available information. The
conceptual objective is to minimize
\begin{equation}
  J(\pi)=\mathbb{E}_{\pi}\!\left[
    \sum_{t=1}^{T}\gamma^{t-1}
    \left\{\ell(\widehat y_t,Y_t)
      +\mathbf{1}_{\{a_t=D\}}C_D(S_t,x_t)\right\}
  \right].
\end{equation}
The expectation is over observations, targets, and any randomness in
prediction, routing, and updating. Their probabilistic specification
remains open. Time and the remaining horizon can be included in
$\mathcal{I}_t$.

Define the current conditional predictive risks by
\begin{align}
  r_F(S_t,x_t)&=\mathbb{E}[\ell(M_F(x_t;\theta_t),Y_t)\mid S_t,x_t],\\
  r_D(S_t,x_t)&=\mathbb{E}[\ell(M_D(x_t),Y_t)\mid S_t,x_t].
\end{align}
Let $V^*(S_{t+1},x_{t+1})$ denote the optimal remaining expected objective
from the next observation, with continuation value zero after $T$.
For a specified probabilistic model and state update, action values can
be written formally as
\begin{align}
  Q_F(S_t,x_t)&=r_F(S_t,x_t)
    +\gamma\mathbb{E}[V^*(S_{t+1},x_{t+1})\mid S_t,x_t,a_t=F],\\
  Q_D(S_t,x_t)&=r_D(S_t,x_t)+C_D(S_t,x_t)\notag\\
    &\quad+\gamma\mathbb{E}[V^*(S_{t+1},x_{t+1})\mid S_t,x_t,a_t=D].
\end{align}
These expressions specify an objective, not a tractable solution or an
algorithm. The optimal choice is $D$ when
\begin{equation}
  Q_D(S_t,x_t)<Q_F(S_t,x_t).
\end{equation}
Ties do not impose a preference between the two actions.

\subsection{Conceptual predictive and learning value}

The current expected risk reduction is
\begin{equation}
  \Delta_{\mathrm{pred}}(S_t,x_t)=r_F(S_t,x_t)-r_D(S_t,x_t).
\end{equation}
Let $V_{\mathrm{future}}(S)$ denote expected future cumulative loss/cost,
with lower values preferred. For a common future horizon, input law, and
continuation criterion, define conceptually
\begin{align}
  \Delta_{\mathrm{learn}}(S_t,x_t)
    &=V_{\mathrm{future}}(S_t)\notag\\
    &\quad-\mathbb{E}_{Z_t}\!\left[
      V_{\mathrm{future}}(U(S_t,x_t,Z_t))\mid S_t,x_t\right].
\end{align}
Here $S_t$ is the reference state without the teacher-induced update;
both terms evaluate the same future period, excluding the current loss
and consultation cost. Dependence of the continuation evaluation on
the current time and observation is implicit and must be made explicit
in a model specialization. If the teacher response is deterministic and
known conditional on $(S_t,x_t)$, the expectation over $Z_t$ is degenerate.

This term represents the expected reduction in future cost/risk caused
specifically by using the teacher response to update the cheap predictor.
It is not a computable closed form, and its sign is not fixed. If future
consultation costs are included in $V_{\mathrm{future}}$, it is not solely
a reduction in predictive loss.

The central candidate decomposition emerging from the sequential
formulation is
\begin{equation}
  \Delta_{\mathrm{pred}}(S_t,x_t)
    +\gamma\Delta_{\mathrm{learn}}(S_t,x_t)>C_D(S_t,x_t).
\end{equation}
This is not a theorem. Comparing the general action values gives the
immediate predictive advantage plus the discounted difference of their
continuation terms, against consultation cost. Identifying that
continuation difference with $\Delta_{\mathrm{learn}}$ requires compatible
continuation policies, dynamics, information, and horizons, with the
action-dependent difference represented by $U$. Other action effects or
updates under $F$ must be accounted for in the reference continuation;
otherwise the more general $Q_D<Q_F$ comparison must be retained. The
precise sufficient assumptions and any computational guarantees remain
to be derived.

\subsection{No assumed scalar confidence threshold}

A scalar confidence threshold requires additional structural assumptions.
Neither its existence nor the sign or functional form of a threshold
correction is established. Learning value need not depend only on
epistemic uncertainty: a highly uncertain sample may have little future
value if its region is rarely encountered again, whereas a moderately
uncertain sample may be valuable in a high-probability future region.
These are conceptual possibilities, not demonstrated results for a
specified model.

Thus $\Delta_{\mathrm{learn}}$ may depend on $x_t$, $S_t$ and $\theta_t$,
the future input distribution, teacher behaviour and error structure,
the update operator $U$, the effective future horizon, and other state
variables. No Gaussian uncertainty model or sufficiency of the cheap
predictor's confidence is assumed.

\subsection{Open theoretical questions}

\begin{enumerate}
  \item Under what assumptions can $\Delta_{\mathrm{learn}}$ be computed
    or approximated locally?
  \item When can the optimal policy be reduced to an interpretable
    decision rule?
  \item What is the relationship with value of information, dual control,
    active learning, knowledge distillation, learning to defer, and
    contextual bandits?
  \item Under what conditions is a myopic policy that ignores
    $\Delta_{\mathrm{learn}}$ strictly suboptimal? In particular, can a
    specified model admit $\Delta_{\mathrm{pred}}<C_D$ but
    $\Delta_{\mathrm{pred}}+\gamma\Delta_{\mathrm{learn}}>C_D$?
  \item Can bounds be obtained for approximations of future value?
  \item What minimal family of models permits analytical results without
    turning the problem into an intractable generic Markov decision process?
\end{enumerate}

The connections to myopic cost-aware routing when
$\Delta_{\mathrm{learn}}=0$, and to active acquisition when immediate
predictive benefit is ignored, remain to be formalized within the
sequential model. No novelty claim or theoretical contribution is
established by this candidate formulation.
